% Standalone source; compile with pdfLaTeX (three passes) or latexmk -pdf.
\documentclass[11pt,a4paper]{article}
\usepackage[T1]{fontenc}
\usepackage[utf8]{inputenc}
\usepackage{lmodern}
\usepackage[a4paper,margin=24mm,headheight=15pt]{geometry}
\usepackage{amsmath,amssymb,amsthm,mathtools,mathrsfs}
\usepackage{microtype,booktabs,longtable,array,enumitem}
\usepackage[dvipsnames]{xcolor}
\usepackage{xurl,fancyhdr,listings}
\usepackage[colorlinks=true,linkcolor=MidnightBlue,citecolor=MidnightBlue,
 urlcolor=MidnightBlue,bookmarksnumbered=true]{hyperref}
\usepackage[nameinlink,noabbrev]{cleveref}
\definecolor{Ink}{HTML}{19354D}
\definecolor{Muted}{HTML}{53616D}
\hypersetup{pdftitle={Differential Algebra of Surreal and Surcomplex Numbers},
 pdfauthor={AI-assisted research exposition prepared for Vladimir Reshetnikov},
 pdfsubject={Differential Hahn workspaces, bounded phases, and oscillation obstructions}}
\newtheorem{theorem}{Theorem}[section]
\newtheorem{proposition}[theorem]{Proposition}
\newtheorem{lemma}[theorem]{Lemma}
\newtheorem{corollary}[theorem]{Corollary}
\theoremstyle{definition}
\newtheorem{definition}[theorem]{Definition}
\newtheorem{example}[theorem]{Example}
\newtheorem{remark}[theorem]{Remark}
\newcommand{\No}{\mathbf{No}}
\newcommand{\SC}{\mathbf{No}[i]}
\newcommand{\R}{\mathbb R}
\newcommand{\C}{\mathbb C}
\newcommand{\Q}{\mathbb Q}
\newcommand{\Z}{\mathbb Z}
\newcommand{\N}{\mathbb N}
\newcommand{\D}{\partial}
\newcommand{\OO}{\mathcal O}
\newcommand{\mm}{\mathfrak m}
\newcommand{\BB}{\mathcal B}
\newcommand{\JJ}{\mathcal J}
\newcommand{\HH}{\mathscr H}
\newcommand{\supp}{\operatorname{supp}}
\newcommand{\st}{\operatorname{st}}
\newcommand{\cis}{\operatorname{cis}}
\newcommand{\Expzero}{\operatorname{Exp}_0}
\newcommand{\Logzero}{\operatorname{Log}_0}
\newcommand{\dlog}{\operatorname{dlog}_{\D}}
\newcommand{\Ob}{\operatorname{ob}_{\D}}
\newcommand{\Ip}{\mathsf I_0}
\newcommand{\pinf}{\pi_{\infty}}
\newcommand{\dd}{\mathop{}\!\mathrm d}
\newcommand{\rank}{\operatorname{rank}}
\newcommand{\GL}{\operatorname{GL}}
\newcommand{\Real}{\operatorname{Re}}
\newcommand{\Imag}{\operatorname{Im}}
\newcommand{\Span}{\operatorname{span}}
\newcommand{\Frob}{\operatorname{Frob}}
\DeclareMathOperator{\Arg}{Arg}
\DeclareMathOperator{\Ker}{ker}
\DeclareMathOperator{\im}{im}
\setlist{nosep,leftmargin=1.6em}
\setlength{\parskip}{4pt}
\setlength{\parindent}{1em}
\emergencystretch=2em
\pagestyle{fancy}
\fancyhf{}
\fancyhead[L]{\small\color{Muted}Differential algebra of surreal and surcomplex numbers}
\fancyhead[R]{\small\color{Muted}September 2026}
\fancyfoot[C]{\thepage}
\renewcommand{\headrulewidth}{0.3pt}
\lstset{basicstyle=\small\ttfamily,columns=fullflexible,breaklines=true,
 frame=single,framerule=0.3pt,showstringspaces=false}
\numberwithin{equation}{section}
\newcolumntype{L}[1]{>{\raggedright\arraybackslash}p{#1}}

\begin{document}
\hypersetup{pageanchor=false}
\begin{titlepage}
\centering
\vspace*{12mm}
{\small\sffamily\color{Muted}A DOCUMENTATION GAP STUDY\par}
\vspace{9mm}
{\Huge\bfseries\color{Ink}Differential Algebra of\par
Surreal and Surcomplex\par Numbers\par}
\vspace{7mm}
{\Large Differential Hahn workspaces, bounded phases,\par
logarithmic derivatives, and oscillation obstructions\par}
\vspace{10mm}
{\large Prepared for the \texttt{Surreal} documentation\par}
{\normalsize Vladimir Reshetnikov\par}
\vspace{4mm}
{\normalsize September 21, 2026\par}
\vspace{11mm}
\begin{minipage}{0.91\textwidth}
\begin{abstract}
The documentation of the \texttt{Surreal} repository develops an extensive
Hahn-supported theory of surcomplex functions, but deliberately separates
that theory from derivations of the scalar field. This article supplies a
bridge using the Berarducci--Mantova derivation $\D$, normalized by
$\D\omega=1$.

The first bridge is support-theoretic: every set of surcomplex numbers lies
in a set-sized, algebraically closed Hahn field stable under $\D$. A
countable closure of its exponent group constructs such a workspace, and
strong additivity carries the construction to common-domain analytic
coefficient families. The second bridge is differential-algebraic. Writing
$\OO$ for the finite real surreals, we prove the exact identity
\[
 \dlog(\SC^\times)=\No+i\D\OO.
\]
Consequently $\D y=(a+ib)y$ has a nonzero surcomplex solution exactly when
$b$ has a finite surreal primitive. An explicit, purely infinite phase
obstruction classifies all rank-one equations up to multiplicative gauge.
For rational $b(\omega)$ the criterion reduces to a degree comparison.
Constant-coefficient equations admit a complete solution-space description:
only their real characteristic roots contribute. In particular,
$\D^2y+y=0$ has no nonzero solution in $\SC$, despite algebraic closedness
and surjectivity of $\D$.

The deep existence theorem for $\D$ is imported with attribution; the
workspace construction, phase analysis, classifications, and coefficient
interfaces are proved here. An explicit Picard--Vessiot extension explains
where an oscillatory solution can live. Exact symbolic checks accompany the
exposition, but are not a formal verification or a claim of priority.
\end{abstract}
\end{minipage}
\vfill
{\small\color{Muted}Repository snapshot:
\texttt{e260237db9b71da8b74a0c13c8e6355119091100}\par}
\end{titlepage}
\pagenumbering{roman}
\hypersetup{pageanchor=true}
\tableofcontents
\clearpage
\pagenumbering{arabic}

\section{The gap and the scope of this article}\label{sec:gap}

\subsection{What the repository already covers}
This article reviews \path{VladimirReshetnikov/Surreal} at commit
\begin{quote}\small\ttfamily
e260237db9b71da8b74a0c13c8e6355119091100
\end{quote}
Its documentation map lists fifteen
research packages in three families: surreal arithmetic and constructions,
surcomplex analysis and geometry, and foundations and computation
\cite{RepoMap}. The topical review for this article used that map, the
catalogue, and the specialist READMEs for analysis, trigonometry,
foundations, and computer algebra. It was not a line-by-line verification
of the collection or an exhaustive absence search through every source
manuscript.

The gap is particularly well documented by the collection itself. The
trigonometry report explicitly leaves the choice of a derivation on $\No$
outside its scope. The foundations report insists that a scalar-field
derivation and a derivative of a function must remain distinct interfaces.
The analysis report develops common-domain holomorphic coefficient
families rather than a selected differential-field structure
\cite{RepoTrig,RepoFound,RepoAnalysis}. Thus the missing topic is not
``calculus'' in general. It is the \emph{interaction of a specified surreal
field derivation with the existing analytic and phase structures}.

A natural home for this article is
\path{docs/surcomplex/differential-algebra/}. It complements the analysis,
trigonometry, and foundations reports; it does not replace them. In
particular, it preserves the distinction between fixed-domain holomorphic
families, arbitrary formal germs, and pointwise derivatives.

\subsection{Four operations that must not be conflated}
Fix $t=\omega^{-1}$. The following comparison already prevents a serious
category error.
\begin{center}
\begin{tabular}{L{0.25\textwidth}L{0.39\textwidth}L{0.25\textwidth}}
\toprule
Operation & What varies & Basic example\\
\midrule
Field derivation $\D$ & The scalar's differential-algebraic meaning &
$\D\omega=1$, $\D t=-t^2$\\
Formal derivative $\partial_Z$ & An indeterminate; coefficients are fixed &
$\partial_Z(\omega Z)=\omega$\\
Coefficient derivation $\D_H$ & Hahn scalar coefficients; the external variable is fixed &
$\D_H(\omega z)=z$\\
Fine pointwise derivative & A function's argument in its stated topology &
The constant function $Z\mapsto\omega$ has derivative $0$\\
\bottomrule
\end{tabular}
\end{center}
A function derivative does not act on its constant coefficients. A field
derivation does. A chain rule relating the two needs an evaluation theorem;
it is not a consequence of using the same prime symbol.

\subsection{The principal results and their status}
The support construction in \cref{sec:hull} makes the scalar derivation
available in a genuine set-sized Hahn workspace. The phase construction in
\cref{sec:phase} identifies the logarithmic derivative of a unit-modulus
surcomplex number. These combine in \cref{sec:logder} to yield an exact
criterion for first-order homogeneous equations and a canonical gauge
classification. \Cref{sec:constant} classifies constant-coefficient linear
equations of every finite order. \Cref{sec:analytic} then returns to the
repository's common-domain coefficient calculus and proves the relevant
commutation and evaluation formulas.

These are expository developments from explicitly stated hypotheses, not a
claim that the individual consequences are new to the literature. The
construction and surjectivity of the Berarducci--Mantova derivation are deep
published results, not proved anew here. The paper does not claim a solution
to a named open problem, a universal differential-equation solver, or a
machine-checked formalization.

\section{Set-supported normal forms and finite parts}\label{sec:support}

\subsection{Class notation and actual sets}
We use classical set theory with definable classes, or an appropriate
conservative class-theoretic presentation. The symbols $\No$ and $\SC$
denote proper classes. A ``class field'' assertion means that the field
operations are class functions satisfying the usual identities. Every
individual number, every support, every family to be summed, and every
polynomial used below is a set. No proper-class sum is taken.

Write a surreal normal form as
\begin{equation}\label{eq:hahn}
 x=\sum_{\gamma\in S}a_\gamma t^\gamma,
 \qquad t^\gamma:=\omega^{-\gamma},\qquad a_\gamma\in\R,
\end{equation}
where $S$ is a well-ordered set of surreal exponents. For surcomplex numbers
the coefficients belong to $\C$. Here $t^\gamma$ is \emph{notation for the
Conway monomial} $\omega^{-\gamma}$. It is not a declaration that an
arbitrary surreal exponent has been interpreted through the exponential
function. Normal forms and strong summation are the background series
structure used in \cite{BM,ADH}.

If $\Gamma\subseteq\No$ is a set-sized additive subgroup, define
\[
 H_\Gamma=\R((t^\Gamma)),\qquad
 K_\Gamma=\C((t^\Gamma))=H_\Gamma[i].
\]
These are actual sets, canonically embedded in $\No$ and $\SC$. Their
supports may be uncountable, but are always subsets of the set $\Gamma$.
When $\Gamma$ is divisible, $H_\Gamma$ is real closed and $K_\Gamma$ is
algebraically closed. We use this standard Hahn-field fact; it concerns
\emph{algebraic} equations, not differential equations \cite{ADH}.

\begin{definition}[Strong summability]\label{def:summable}
A set-indexed family $(x_j)_{j\in J}$ of Hahn series is summable if
$\bigcup_j\supp x_j$ is well ordered and, for each exponent $\gamma$, only
finitely many $j$ have a nonzero coefficient at $\gamma$. Its sum is formed
by those finite coefficient sums.
\end{definition}
This definition allows coefficientwise rearrangement when the combined
family is summable. It does not authorize an arbitrary exchange of two
infinite sums. In particular, writing a formal geometric or exponential
series is not by itself a summability proof.

We use the positive-support form of the Hahn--Neumann lemma: the additive
monoid generated by a well-ordered subset of positive exponents is well
ordered, and each exponent has only finitely many representations by a
finite nondecreasing list from that subset. Consequently formal power
series with ordinary real or complex coefficients can be evaluated at
infinitesimals. This is the ordered-series lemma behind the standard
Hahn calculus; see \cite{Neumann} and the summability preliminaries of
\cite{BM}. Every use below has positive infinitesimal support.

\subsection{The valuation ring and the purely infinite projection}
For real surreals put
\[
 \OO=\{x: |x|\le r\text{ for some }r\in\R_{>0}\},\qquad
 \mm=\{x: |x|<1/n\text{ for every }n\ge1\}.
\]
Thus $\OO$ is the ring of finite elements and $\mm$ its maximal ideal.
Let $\JJ$ be the additive class of purely infinite normal forms: all
nonzero terms have $\gamma<0$ in \eqref{eq:hahn}. Splitting a support at
zero gives the unique additive decomposition
\begin{equation}\label{eq:threeparts}
 \No=\JJ\oplus\R\oplus\mm,
 \qquad x=x_\infty+c_0(x)+x_\mathrm{small}.
\end{equation}
The maps $\pinf(x)=x_\infty$ and $c_0(x)$ are $\R$-linear. Neither is
asserted to be multiplicative. On $\OO$, $c_0$ is the standard-part map.

For $z=a+ib$, let $\bar z=a-ib$ and
$|z|=(a^2+b^2)^{1/2}$. Complex finite and infinitesimal elements form
$\OO_\C=\OO+i\OO$ and $\mm_\C=\mm+i\mm$. Their standard parts lie in
$\C$. A unit-modulus element is necessarily finite, irrespective of the
birthdays of its real and imaginary parts.

\begin{remark}[A quotient-class safeguard]\label{rem:quotient}
The suggestive notation $\No/\OO$ must not be implemented as a collection
whose members are proper-class cosets. In this paper it always means the
\emph{represented quotient} $\JJ$, with quotient map $\pinf$. Its
representatives are surreal numbers, hence sets. All later obstruction
classes will likewise be given by actual surreal representatives. In a
fixed set-sized field, ordinary set quotients are of course legitimate.
\end{remark}

\section{The selected scalar derivation}\label{sec:derivation}

\subsection{The imported existence theorem}
\begin{theorem}[Berarducci--Mantova, imported]\label{thm:BM}
There is a specified derivation $\D:\No\to\No$ satisfying
\begin{align*}
 \D(xy)&=x\D y+y\D x, & \Ker\D&=\R,\\
 \D\!\left(\sum_j x_j\right)&=\sum_j\D x_j,
 & \D(e^x)&=e^x\D x,
\end{align*}
where strong additivity includes summability of the differentiated family
whenever the original family is summable. Moreover,
\[
 x>\R\Longrightarrow\D x>0,\qquad
 \D\omega=1,\qquad \D\No=\No.
\]
Here $e^x$ is the Gonshor exponential. We use the particular
Berarducci--Mantova derivation with these properties.
\end{theorem}
The existence and integration assertions are Theorems A and B of
\cite{BM}. They are assumptions imported from that paper, not consequences
of the elementary calculations that follow. The normalization
$\D\omega=1$ alone is not being claimed to characterize a derivation on
all of $\No$.

The expression ``Liouville closed'' in the real differential-field setting
includes real closedness, primitives, and exponential integrals. In our
setting the latter follow directly: if $\D A=a$, then
\begin{equation}\label{eq:realexpintegral}
 \frac{\D(e^A)}{e^A}=a.
\end{equation}
This observation will fail for a general \emph{complex} coefficient, not
because primitives fail, but because a compatible global complex
exponential is unavailable.

\subsection{Elementary differential identities}
\begin{proposition}\label{prop:calculus}
For $x\ne0$, $u>0$, and $r\in\R$,
\[
 \D(x^{-1})=-x^{-2}\D x,\qquad
 \D\log u=\frac{\D u}{u},\qquad
 \D(u^r)=r u^{r-1}\D u,
\]
where $u^r=\exp(r\log u)$. In particular
\[
 \D\omega^r=r\omega^{r-1},\qquad
 \D t^r=-r t^{r+1}\quad(r\in\R).
\]
\end{proposition}
\begin{proof}
Differentiate $xx^{-1}=1$ and $\exp(\log u)=u$. Then apply exponential
compatibility to $\exp(r\log u)$, using $\D r=0$. For ordinary real $r$,
the usual real-power interpretation of $\omega^r$ agrees with its Conway
monomial interpretation; this is the real-exponent special case. Finally
$t=\omega^{-1}$.
\end{proof}
For an arbitrary surreal $\gamma$, one must \emph{not} infer
$\D(\omega^\gamma)=\omega^\gamma\D(\gamma\log\omega)$ by identifying the
omega-map with generic exponentiation. The safe universal identity is
\[
 \D(\omega^\gamma)=\omega^\gamma\D\log(\omega^\gamma).
\]
The latter logarithm is the actual Gonshor logarithm of the monomial. The
workspace theorem below intentionally uses its derivative as supplied by
$\D$, rather than an incorrect universal exponent rule.

\begin{proposition}[Complexification]\label{prop:complexification}
There is a unique derivation extending $\D$ to $\SC$, namely
\[
 \D(a+ib)=\D a+i\D b.
\]
It commutes with conjugation, is strongly additive, has constant field
$\C$, and is surjective.
\end{proposition}
\begin{proof}
Differentiating $i^2+1=0$ gives $2i\D i=0$, hence $\D i=0$. This forces
the displayed formula, and expansion of products verifies Leibniz's rule.
The kernel and surjectivity follow componentwise from \cref{thm:BM}.
Strong summability and strong additivity pass through the real and
imaginary parts, since only finite unions of supports are involved.
\end{proof}

\begin{proposition}[Algebraic implicit differentiation]\label{prop:algdiff}
Let $F$ be a differential subfield of $\SC$, and let $\alpha$ be algebraic
over $F$. If $P(X)\in F[X]$ is its minimal polynomial, then
\[
 \D\alpha=-\frac{P^\D(\alpha)}{P_X(\alpha)},\qquad
 P^\D(X)=\sum_k(\D a_k)X^k\quad\text{for }P=\sum_k a_kX^k.
\]
\end{proposition}
\begin{proof}
In characteristic zero the irreducible polynomial $P$ is separable, so
$P_X(\alpha)\ne0$. Differentiating the finite identity $P(\alpha)=0$ gives
$P^\D(\alpha)+P_X(\alpha)\D\alpha=0$.
\end{proof}
This is the appropriate derivative rule for an exact algebraic root
object. It is not a choice of a numerical root branch.

\subsection{A normalized primitive without an arbitrary constant}
\begin{proposition}\label{prop:normalizedprimitive}
For each $a\in\No$ there is a unique $\Ip(a)\in\No$ with
\[
 \D\Ip(a)=a,\qquad c_0(\Ip(a))=0.
\]
The map $\Ip$ is $\R$-linear and
$\Ip(\D x)=x-c_0(x)$.
\end{proposition}
\begin{proof}
Take any primitive and subtract its real coefficient at exponent zero.
Any two primitives differ by a real constant, proving both existence and
uniqueness under the normalization. Linearity and the last identity follow
from uniqueness.
\end{proof}
The definition is a legitimate class function by uniqueness. It is not an
algorithm for computing arbitrary primitives, and it does not assert
strong additivity of integration.

\section{Differentially stable Hahn workspaces}\label{sec:hull}

The repository's set-sized workspace principle is algebraic. To use a
scalar derivation on an entire Hahn workspace, its exponent group must
also accommodate the supports created by differentiating monomials.
This section supplies that missing closure condition.

\begin{definition}
A set-sized additive subgroup $\Gamma\subseteq\No$ is
\emph{differentially closed on monomials} if
\[
 \supp\D(t^\gamma)\subseteq\Gamma\qquad(\gamma\in\Gamma).
\]
The support here and below is in the $t$-exponent convention of
\eqref{eq:hahn}.
\end{definition}

\begin{lemma}[Monomial criterion]\label{lem:monomialcriterion}
For a set-sized additive group $\Gamma$, the following are equivalent:
\[
 \D K_\Gamma\subseteq K_\Gamma
 \quad\Longleftrightarrow\quad
 \Gamma\text{ is differentially closed on monomials}.
\]
The same equivalence holds for $H_\Gamma$.
\end{lemma}
\begin{proof}
Necessity follows by applying the field-closure condition to each
$t^\gamma$. For sufficiency, write
$x=\sum_{\gamma\in S}c_\gamma t^\gamma$ with $S\subseteq\Gamma$.
Strong additivity gives the summable family
$(c_\gamma\D t^\gamma)_{\gamma\in S}$ and
\[
 \D x=\sum_{\gamma\in S}c_\gamma\D t^\gamma.
\]
Every term has support in $\Gamma$, so the sum does too. Ordinary real
and complex coefficients are constants of the derivation.
\end{proof}

\begin{theorem}[Countable differential Hahn hull]\label{thm:hull}
Given a set-sized additive subgroup $\Gamma_0\subseteq\No$, define
\begin{align}
 \Gamma_{n+1}
 &=\Span_\Q\!\left(\Gamma_n\cup
        \bigcup_{\gamma\in\Gamma_n}\supp\D(t^\gamma)\right),
       \label{eq:hullstage}\\
 \Gamma_*&=\bigcup_{n<\omega}\Gamma_n.
       \label{eq:hullunion}
\end{align}
Then $\Gamma_*$ is a set-sized divisible additive group,
$K_{\Gamma_*}$ is algebraically closed, and
$\D K_{\Gamma_*}\subseteq K_{\Gamma_*}$. It is the least such divisible
monomial-closed exponent group containing $\Gamma_0$.
\end{theorem}
\begin{proof}
Each derivative in \eqref{eq:hullstage} is an individual surreal and
therefore has a set support. Replacement and union show that the displayed
union is a set. The rational span is again a set. The stages are increasing,
so their countable union is a set-sized additive group. Every member lies
in a stage whose next stage contains all its rational multiples; thus the
union is divisible.

If $\gamma\in\Gamma_*$, it belongs to some $\Gamma_n$. The entire support
of $\D(t^\gamma)$ is then in $\Gamma_{n+1}\subseteq\Gamma_*$. Apply
\cref{lem:monomialcriterion} and the divisible-group Hahn-field theorem.
For minimality, any divisible monomial-closed group containing $\Gamma_0$
contains each stage by induction and hence their union.
\end{proof}

\begin{corollary}[Differential localization of arbitrary set data]
Every set $A\subseteq\SC$ lies in a set-sized algebraically closed
Hahn field stable under $\D$.
\end{corollary}
\begin{proof}
Take $\Gamma_0$ to be the additive group generated by the union of the
supports of the members of $A$, and apply \cref{thm:hull}.
\end{proof}
The group need not be countable: it is the \emph{number of closure stages}
that is countable. Also, this is a minimality statement among full Hahn
fields indexed by divisible exponent groups, not among all differential
subfields.

\subsection{An exact finite-rank example}
\begin{proposition}\label{prop:realgroup}
Let $0\ne\Gamma\subseteq\R$ be an additive subgroup. Then $K_\Gamma$ is
stable under $\D$ if and only if $1\in\Gamma$. If $\Gamma$ is divisible,
its differential Hahn hull is $\Gamma+\Q$.
\end{proposition}
\begin{proof}
For nonzero $\gamma\in\Gamma$, \cref{prop:calculus} gives
$\D t^\gamma=-\gamma t^{\gamma+1}$. Closure therefore requires
$\gamma+1\in\Gamma$, hence $1\in\Gamma$. Conversely, containing $1$
provides every exponent $\gamma+1$. The rational-span assertion follows
from \cref{thm:hull}. The exceptional group $\{0\}$ gives the constant
field and is already stable.
\end{proof}
For example, $\Gamma=\Q\sqrt2$ is divisible but not differentially stable:
\[
 \D t^{\sqrt2}=-\sqrt2\,t^{\sqrt2+1}\notin K_{\Q\sqrt2}.
\]
Its hull is $\Q+\Q\sqrt2$. This is a simple counterexample to the claim
that an arbitrary algebraically closed Hahn workspace automatically
supports the ambient derivation.

\subsection{Differentiation closure is not integration closure}
\begin{example}\label{ex:logintegral}
The field $H_\Q=\R((t^\Q))$ is stable under $\D$, but $t$ has no primitive
in it. Its primitives in $\No$ are $\log\omega+c$, $c\in\R$.
\end{example}
\begin{proof}
For $x=\sum c_qt^q$, strong additivity gives
$\D x=\sum(-q)c_qt^{q+1}$. The coefficient at exponent $1$ is always zero:
the only possible source is $q=0$, whose multiplier is zero. The derivative
therefore cannot equal $t$. On the other hand,
$\D\log\omega=\omega^{-1}=t$.
\end{proof}
Thus the whole-field surjectivity of \cref{thm:BM} must not be silently
transported to a localized workspace. The counterexample is about
integration, not about algebraic closedness or Hahn completeness. General
differentially stable Hahn subfields also occur in the model-theoretic
analysis of \cite{ADH}; our elementary construction is designed specifically
to expose the support contract needed here.

\section{Finite phases and their scalar derivatives}\label{sec:phase}

The finite-angle theorem is already central to the repository's
trigonometry report \cite{RepoTrig}. We give the short formal-series route
needed here, then prove its compatibility with the selected scalar
derivation. This extra compatibility is the point of the section.

\subsection{Small exponential and logarithm}
For $h\in\mm_\C$, define the Hahn sums
\begin{equation}\label{eq:smallfunctions}
 \Expzero(h)=\sum_{n\ge0}\frac{h^n}{n!},\qquad
 \Logzero(1+h)=\sum_{n\ge1}\frac{(-1)^{n+1}}n h^n.
\end{equation}
Positive support and the Hahn--Neumann lemma justify both sums. Formal
power-series identities over $\Q$ therefore specialize to
\begin{align}
 \Expzero(h+k)&=\Expzero(h)\Expzero(k),\label{eq:smallgroup}\\
 \Logzero(\Expzero(h))&=h,\qquad
 \Expzero(\Logzero(1+h))=1+h.\label{eq:smallinverse}
\end{align}
To be explicit about the substitution principle, a coefficient at a fixed
Hahn exponent receives contributions from only finitely many finite
products of the supports of $h$ and $k$. The formal identity can thus be
checked at that exponent by finite algebra. This is not an appeal to a
limit of partial sums.

\begin{lemma}[Scalar chain rules for small series]\label{lem:smallchain}
For $h\in\mm_\C$,
\[
 \D\Expzero(h)=\Expzero(h)\D h,\qquad
 \D\Logzero(1+h)=\frac{\D h}{1+h}.
\]
\end{lemma}
\begin{proof}
Strong additivity and the finite power rule give
\[
 \D\sum_{n\ge0}\frac{h^n}{n!}
 =\sum_{n\ge1}\frac{h^{n-1}\D h}{(n-1)!}
 =\D h\,\Expzero(h).
\]
Multiplication by the fixed number $\D h$ preserves summability. The same
argument differentiates the logarithm series and gives
$\D h\sum_{n\ge0}(-h)^n=\D h/(1+h)$. All evaluated series have positive
infinitesimal support before the fixed multiplier is applied.
\end{proof}

\subsection{The entire unit circle has finite phase}
Let
\[
 T(\No)=\{u\in\SC:u\bar u=1\}.
\]
For $\theta=r+\varepsilon\in\OO$, with $r\in\R$ and
$\varepsilon\in\mm$, set
\begin{equation}\label{eq:cis}
 \cis(\theta)=e^{ir}\Expzero(i\varepsilon).
\end{equation}
The factor $e^{ir}$ is the ordinary complex exponential. No value of a
circular function at an infinite angle is used.

\begin{theorem}[Finite-phase parametrization]\label{thm:phase}
The map $\cis:(\OO,+)\to(T(\No),\cdot)$ is a surjective group
homomorphism with kernel $2\pi\Z$. Every $u\in T(\No)$ has a unique
factorization
\begin{equation}\label{eq:normalizedphase}
 u=c\Expzero(i\varepsilon),\qquad
 c\in T(\R),\quad\varepsilon\in\mm.
\end{equation}
Moreover,
\begin{equation}\label{eq:phasechain}
 \frac{\D\cis(\theta)}{\cis(\theta)}=i\D\theta.
\end{equation}
\end{theorem}
\begin{proof}
The group law follows from \eqref{eq:smallgroup}. Conjugation carries
$\Expzero(i\varepsilon)$ to $\Expzero(-i\varepsilon)$, so its norm is
one. Conversely, let $u\bar u=1$. Then $u$ is finite, and
$c=\st(u)\in\C$ satisfies $c\bar c=1$. Write $u/c=1+h$, where
$h\in\mm_\C$, and put $\eta=\Logzero(1+h)$. Since
$\overline{1+h}=(1+h)^{-1}$, the formal logarithm identities give
\[
 \bar\eta=\Logzero((1+h)^{-1})=-\eta.
\]
Thus $\eta=i\varepsilon$ for a real infinitesimal $\varepsilon$, and
$u=c\Expzero(i\varepsilon)$. Choose an ordinary real argument $r$ of
$c$ to obtain $u=\cis(r+\varepsilon)$. The normalized factorization is
unique because $c=\st(u)$ and $\Logzero$ is inverse to $\Expzero$ on
infinitesimals.

If $\cis(r+\varepsilon)=1$, standard parts give $e^{ir}=1$, so
$r\in2\pi\Z$. Then \eqref{eq:smallinverse} gives $\varepsilon=0$.
Finally, $r$ and $e^{ir}$ are differential constants, and
\cref{lem:smallchain} proves \eqref{eq:phasechain}.
\end{proof}
The angular branch ambiguity is an ordinary constant $2\pi k$, so it
vanishes under $\D$. More strongly, the infinitesimal angle in
\eqref{eq:normalizedphase} is branch-free.

\begin{corollary}[Differential polar decomposition]\label{cor:polar}
Every nonzero $z\in\SC$ can be written as
$z=\rho\cis\theta$, with $\rho=|z|>0$ and $\theta\in\OO$, and
\begin{equation}\label{eq:polarlogder}
 \frac{\D z}{z}=\frac{\D\rho}{\rho}+i\D\theta.
\end{equation}
In particular,
\[
 \Real\frac{\D z}{z}=\frac{\D|z|}{|z|},\qquad
 \Imag\frac{\D z}{z}=\D\theta.
\]
\end{corollary}
\begin{proof}
Apply \cref{thm:phase} to $z/|z|$, and use the product rule. Both terms
on the right side of \eqref{eq:polarlogder} have the indicated real or
purely imaginary type.
\end{proof}

\section{The subgroup of finite-phase speeds}\label{sec:bounded}

\begin{definition}\label{def:B}
The finite-primitive subgroup is
\[
 \BB:=\D\OO=\{\D\theta:\theta\in\OO\}\subseteq\No.
\]
Membership means existence of a \emph{finite} primitive, not merely
existence of some primitive.
\end{definition}

\subsection{Order-theoretic structure}
\begin{lemma}[Every finite element has infinitesimal derivative]
\label{lem:smallness}
If $f\in\OO$, then $\D f\in\mm$.
\end{lemma}
\begin{proof}
For every positive integer $n$, both $\omega/n+f$ and $\omega/n-f$
are positive infinite. The $H$-field sign condition gives
\[
 1/n+\D f>0,\qquad1/n-\D f>0.
\]
Hence $|\D f|<1/n$ for every $n$.
\end{proof}
This proof uses the normalization $\D\omega=1$ and the sign condition,
not a topological limiting argument.

\begin{theorem}[Structure of $\BB$]\label{thm:Bstructure}
The following assertions hold.
\begin{enumerate}[label=\textup{(\roman*)}]
\item $\BB=\D\mm\subsetneq\mm$.
\item $\D:\mm\to\BB$ is an $\R$-linear bijection that reverses strict
order.
\item $\BB$ is a convex additive subgroup of $\No$ and an $\OO$-submodule
of the additive group of $\No$.
\item An element $b$ belongs to $\BB$ exactly when $\Ip(b)$ is finite,
equivalently when $\pinf\Ip(b)=0$.
\end{enumerate}
\end{theorem}
\begin{proof}
Since $\OO=\R\oplus\mm$ and $\D\R=0$, the first equality holds.
\Cref{lem:smallness} gives containment in $\mm$. The kernel of the
restriction to $\mm$ is zero, because a real infinitesimal constant is
zero. If $0<\varepsilon\in\mm$, then $1/\varepsilon>\R$, so
\[
 0<\D(1/\varepsilon)=-\varepsilon^{-2}\D\varepsilon.
\]
Thus $\D\varepsilon<0$, proving the order reversal.

For convexity, suppose $0<b\le c$ with $c\in\BB$. Choose finite
$\theta$ with $\D\theta=c$, and choose a primitive $u$ of $b$. If $u$
were negative infinite, its derivative would be negative by the sign
condition. If $u$ were positive infinite, $u-\theta$ would be positive
infinite, giving $b-c>0$, a contradiction. Therefore $u$ is finite and
$b\in\BB$. The zero case and negative elements follow by additivity.
For $a\in\OO$ choose $r\in\R_{>0}$ with $|a|\le r$. If $b\in\BB$,
then $|ab|\le r|b|$, and convexity together with real linearity gives
$ab\in\BB$.

All primitives of $b$ differ by real constants, so either all are finite
or none are. This proves (iv). Finally $t=1/\omega$ is infinitesimal but
has primitive $\log\omega$, which is positive infinite. Hence
$t\notin\BB$, proving strictness. In contrast,
$t^2=\D(-t)\in\BB$.
\end{proof}

\begin{proposition}[A represented differential quotient]\label{prop:decompB}
There is a direct-sum decomposition of additive $\R$-vector classes
\begin{equation}\label{eq:Bsplit}
 \No=\D\JJ\oplus\BB.
\end{equation}
The projection onto the first summand is
$b\mapsto\D\pinf\Ip(b)$. The restriction
$\D:\JJ\to\D\JJ$ is an order-preserving $\R$-linear bijection.
\end{proposition}
\begin{proof}
Decompose a normalized primitive as
$\Ip(b)=j+\varepsilon$, with $j\in\JJ$, $\varepsilon\in\mm$.
Then $b=\D j+\D\varepsilon$. If $\D j\in\BB$, choose finite
$\theta$ with $\D j=\D\theta$. The difference $j-\theta$ is real,
forcing the purely infinite $j$ to be finite, hence zero. This proves the
direct sum. A nonzero positive member of $\JJ$ is positive infinite,
so its derivative is positive. Apply this to differences to obtain the
order assertion.
\end{proof}
This also explains the represented version of the familiar quotient
isomorphism $\No/\OO\cong\No/\BB$ induced by differentiation. The actual
representatives are $\JJ$ and $\D\JJ$; no proper-class coset is made into
an object.

\subsection{A hierarchy of explicit thresholds}
For real $p$ the equation $\D\theta=\omega^{-p}$ has the primitives
\begin{equation}\label{eq:powerprimitive}
 \theta=
 \begin{cases}
 \displaystyle\frac{\omega^{1-p}}{1-p}+c,&p\ne1,\\[5pt]
 \log\omega+c,&p=1,
 \end{cases}
 \qquad c\in\R.
\end{equation}
Consequently
\begin{equation}\label{eq:powerthreshold}
 \omega^{-p}\in\BB\quad\Longleftrightarrow\quad p>1.
\end{equation}
Thus even a positive infinitesimal frequency can be inadmissible.

The same phenomenon repeats at each finite logarithmic depth. Put
$L_0=\omega$ and $L_{k+1}=\log L_k$. For $k\ge0$, with an empty product
interpreted as $1$, set
\[
 b_{k,p}=\frac{1}{L_0L_1\cdots L_{k-1}L_k^p}.
\]
The chain rule yields
\[
 \D L_k=\frac1{L_0\cdots L_{k-1}},\qquad
 \Ip(b_{k,p})=
 \begin{cases}
 L_k^{1-p}/(1-p),&p\ne1,\\
 L_{k+1},&p=1.
 \end{cases}
\]
All these primitives have zero ordinary constant coefficient. Each $L_k$
is positive infinite, so $b_{k,p}\in\BB$ exactly when $p>1$.
In particular, neither $1/\omega$ nor
$1/(\omega\log\omega)$ belongs to $\BB$; both are infinitesimal.

\section{Logarithmic derivatives and complete rank-one classification}
\label{sec:logder}

For $y\ne0$ define $\dlog(y)=\D y/y$. It is a homomorphism from the
multiplicative class group to the additive one, since
$\dlog(yz)=\dlog(y)+\dlog(z)$.

\begin{theorem}[Exact logarithmic-derivative image]\label{thm:logimage}
For the selected derivation on $\SC$,
\begin{equation}\label{eq:mainimage}
 \boxed{\quad\dlog(\SC^\times)=\No+i\BB
             =\No+i\D\OO.\quad}
\end{equation}
The kernel of $\dlog$ is $\C^\times$.
\end{theorem}
\begin{proof}
For any nonzero $y$, \cref{cor:polar} gives
$\dlog y=\D\rho/\rho+i\D\theta$, with $\theta$ finite. Thus its
imaginary part lies in $\BB$.

Conversely, let $a\in\No$ and $b\in\BB$. Choose $A=\Ip(a)$ and the
unique $\varepsilon\in\mm$ with $\D\varepsilon=b$. Then
\begin{equation}\label{eq:generator}
 h=\exp(A)\Expzero(i\varepsilon)\ne0
\end{equation}
satisfies $\D h/h=a+ib$ by exponential compatibility and
\cref{lem:smallchain}. Finally $\dlog y=0$ exactly when
$\D y=0$, equivalently $y\in\C^\times$.
\end{proof}

\begin{corollary}[First-order homogeneous equation]\label{cor:rankone}
For $q=a+ib\in\SC$, the following are equivalent:
\begin{enumerate}[label=\textup{(\roman*)}]
\item $\D y=qy$ has a nonzero solution in $\SC$;
\item $b\in\BB$;
\item $\pinf\Ip(b)=0$.
\end{enumerate}
When they hold, every solution is $ch$ with $c\in\C$ and $h$ as in
\eqref{eq:generator}; when they fail, the only solution is zero.
\end{corollary}
\begin{proof}
Only the assertion about all solutions remains. If $h\ne0$ is one
solution, $\D(y/h)=0$, so $y/h$ is an ordinary complex constant.
\end{proof}
In particular, a one-dimensional linear equation need not have the
expected one-dimensional nonzero solution space \emph{inside this field}.
The familiar dimension bound is an upper bound until a sufficiently large
differential extension is supplied.

\subsection{A canonical purely infinite phase obstruction}
\begin{definition}\label{def:obstruction}
For $q\in\SC$ put
\begin{equation}\label{eq:obstruction}
 \Ob(q)=\pinf\Ip(\Imag q)\in\JJ.
\end{equation}
\end{definition}
The notation has an actual surreal representative, not an ambiguous
choice of an antiderivative. The normalization of $\Ip$ is convenient;
subtracting any other real integration constant would give the same
purely infinite part.

\begin{theorem}[Exact obstruction sequence]\label{thm:obstruction}
The map $\Ob:\SC\to\JJ$ is a surjective $\R$-linear map with
\[
 \Ker\Ob=\dlog(\SC^\times).
\]
More explicitly, there is an exact sequence of class groups
\begin{equation}\label{eq:exactsequence}
 1\longrightarrow\C^\times\longrightarrow\SC^\times
 \xrightarrow{\,\dlog\,}(\SC,+)
 \xrightarrow{\,\Ob\,}(\JJ,+)\longrightarrow0.
\end{equation}
Every map in the sequence acts on set-sized elements.
\end{theorem}
\begin{proof}
Linearity follows from \cref{prop:normalizedprimitive} and
\eqref{eq:threeparts}. The kernel is \cref{thm:logimage}. Given
$j\in\JJ$, the normalized primitive of $\D j$ is $j$, so
$\Ob(i\D j)=j$, proving surjectivity. Exactness at the multiplicative
terms is the kernel statement of \cref{thm:logimage}.
\end{proof}

\begin{definition}[Multiplicative gauge equivalence]
Two coefficients $q,r\in\SC$ are gauge equivalent if
$q-r=\dlog u$ for some $u\in\SC^\times$. Equivalently, substitution
$y=uv$ transforms $\D y=qy$ into $\D v=rv$.
\end{definition}

\begin{theorem}[Canonical rank-one representatives]\label{thm:gauge}
Two rank-one equations are gauge equivalent if and only if their
obstructions agree. Each class has a unique representative of the form
\begin{equation}\label{eq:canonicalgauge}
 \D y=i(\D j)y,\qquad j\in\JJ.
\end{equation}
Tensor product adds the representatives' parameters $j$, and passage to
the dual negates $j$.
\end{theorem}
\begin{proof}
The equivalence criterion follows from the kernel in
\cref{thm:obstruction}. For an explicit reduction, write
$\Ip(b)=j+\varepsilon$ with $\varepsilon\in\mm$, and put
$u=\exp(\Ip(a))\Expzero(i\varepsilon)$. Then
\[
 q-\dlog u=i\D j.
\]
The obstruction of the latter coefficient is $j$, proving uniqueness.
The coefficient of a tensor product is the sum of the two scalar
coefficients, and that of a dual is its negative, giving the last claim.
\end{proof}
Examples of the parameter $j$ are informative:
\[
 \Ob(i)=\omega,\qquad
 \Ob(i/\omega)=\log\omega,\qquad
 \Ob(i/\omega^2)=0.
\]
The obstruction measures an \emph{accumulated infinite phase}, not merely
the size of the instantaneous coefficient. A real coefficient has zero
obstruction, however large it is.

\subsection{What happens to inhomogeneous equations}
\begin{proposition}[Variation of constants on the admissible locus]
\label{prop:inhom}
Suppose $\Ob(q)=0$, and let $h\ne0$ satisfy $\D h=qh$. For every
$f\in\SC$, the solutions of
$\D y-qy=f$ are
\[
 y=h(v+c),\qquad\D v=f/h,\qquad c\in\C.
\]
In particular the inhomogeneous equation is solvable.
\end{proposition}
\begin{proof}
Surjectivity of $\D$ on $\SC$ supplies $v$. The product rule verifies
the formula, and subtraction of two solutions reduces to the homogeneous
equation.
\end{proof}
Outside the admissible locus, failure of a nonzero homogeneous solution
does \emph{not} imply failure of every inhomogeneous equation. For example,
\[
 \D y-iy=1
\]
has the unique solution $y=i$: it is a solution by direct substitution,
and any difference of two solutions would solve $\D u=iu$, whose only
solution is zero. In general we make no inhomogeneous solvability claim
outside the locus of \cref{prop:inhom}.

\section{An exact decision criterion for rational coefficients}
\label{sec:rational}

The full obstruction map is a semantic classification rather than a
universal effective procedure. For rational functions over the ordinary
constants, however, the condition is decidable by finite algebra.

\begin{theorem}[Rational phase criterion]\label{thm:rational}
Let $b(X)\in\R(X)$ be nonzero, and write it in reduced form
$b=P/Q$. Then
\begin{equation}\label{eq:rationalcriterion}
 b(\omega)\in\BB
 \quad\Longleftrightarrow\quad
 \deg Q-\deg P\ge2.
\end{equation}
For $q(\omega)=a(\omega)+ib(\omega)$ with $a,b\in\R(X)$, this is
exactly the criterion for a nonzero solution of $\D y=q(\omega)y$.
The case $b=0$ is always admissible.
\end{theorem}
\begin{proof}
Let $p=\deg Q-\deg P$. Expanding the rational function at infinity gives
\[
 b(\omega)=c t^p+\text{higher integer powers of }t,
 \qquad c\in\R^\times.
\]
If $p\ge2$, write $b(\omega)=\sum_{n\ge2}b_nt^n$. The Hahn series
\[
 \theta=-\sum_{n\ge2}\frac{b_n}{n-1}t^{n-1}
\]
is infinitesimal and differentiates strongly to $b(\omega)$. Thus
$b(\omega)\in\BB$.

If $p\le1$, its leading term gives
$|b(\omega)|\ge |c|t^p/2$. Were $b(\omega)$ in $\BB$, real linearity
and convexity would force $t^p\in\BB$. This contradicts
\eqref{eq:powerthreshold}. The assertion for $q$ is
\cref{cor:rankone}.
\end{proof}
No algorithm for an arbitrary real coefficient is asserted: the theorem is
algebraic over $\R$, while the accompanying executable restricts its input
to exact rational coefficients so that polynomial arithmetic and degree
comparison are effective.

\begin{center}
\begin{tabular}{lll}
\toprule
$b(\omega)$ & Finite primitive? & Reason or primitive\\
\midrule
$0$ & yes & $0$\\
$1$ & no & $\omega$ is infinite\\
$1/\omega$ & no & $\log\omega$ is infinite\\
$1/\omega^2$ & yes & $-1/\omega$\\
$\omega/(1+\omega^2)$ & no & degree gap $1$\\
$1/(1+\omega^2)$ & yes & degree gap $2$\\
$(1-\omega^2)/(1+\omega^2)^2$ & yes & $\omega/(1+\omega^2)$\\
\bottomrule
\end{tabular}
\end{center}

\begin{example}[An algebraic solution with a small angular speed]
\label{ex:algebraicphase}
For $b=1/(1+\omega^2)$ the unit-modulus algebraic number
\[
 y=\frac{1+i\omega}{\sqrt{1+\omega^2}}
\]
satisfies $\D y=iby$. Indeed,
\[
 \frac{\D y}{y}
 =\frac{i}{1+i\omega}-\frac{\omega}{1+\omega^2}
 =\frac{i}{1+\omega^2}.
\]
This example needs neither a named inverse trigonometric function nor an
infinite-phase exponential.
\end{example}

\begin{example}[A genuinely infinite Hahn phase solution]
For $q=i/\omega^2$, the normalized solution is
\[
 y=\Expzero(-it)
 =1-it-\frac{t^2}{2}+\frac{it^3}{6}+\frac{t^4}{24}-\cdots.
\]
Its support is contained in $\N$, so the series is strongly summable.
The identity $\D y=it^2y$ is exact, not an asymptotic approximation.
Finite truncations used in the verifier have a computable residual, but
are not the entire solution.
\end{example}

\section{Oscillation obstructions and constant-coefficient equations}
\label{sec:constant}

\subsection{A short independent proof of the basic obstruction}
\begin{theorem}[No scalar harmonic oscillation]\label{thm:noosc}
The only solution in $\SC$ of $\D y=iy$ is zero. More generally,
$\D y=i\beta y$ has no nonzero solution for any
$\beta\in\R\setminus\{0\}$.
\end{theorem}
\begin{proof}
Suppose first that $\D y=iy$ and $y\ne0$. Conjugation gives
$\D\bar y=-i\bar y$, hence
$\D(y\bar y)=0$. Its positive value is a real constant. Dividing $y$ by
its positive real modulus therefore gives $u+iv$ with
\[
 u^2+v^2=1,\qquad \D u=-v,\qquad\D v=u.
\]
Both $u$ and $v$ are finite. By \cref{lem:smallness}, both derivatives
are infinitesimal, and the differential equations then make $u,v$
infinitesimal. Their squared sum cannot be $1$.

For general $\beta$, the same argument gives
$\D u=-\beta v$ and $\D v=\beta u$. Division by the nonzero real
constant $\beta$ again makes both coordinates infinitesimal.
\end{proof}
This proof does not use the full finite-phase parametrization or
surjectivity of $\D$. It uses the constant field, the $H$-field sign
condition, and $\D\omega=1$. The obstruction is therefore not an artifact
of a particular logarithm branch.

\begin{corollary}\label{cor:constantlambda}
For $\lambda\in\C$, the equation $\D y=\lambda y$ has a nonzero
solution in $\SC$ exactly when $\lambda\in\R$. In that case its
solutions are $c e^{\lambda\omega}$, $c\in\C$.
\end{corollary}
\begin{proof}
Write $\lambda=\alpha+i\beta$. Multiplication by
$e^{-\alpha\omega}$ transforms a solution into one of
$\D v=i\beta v$. Apply \cref{thm:noosc}; when $\beta=0$, the quotient
by $e^{\alpha\omega}$ has derivative zero.
\end{proof}

\subsection{All finite-order constant-coefficient equations}
\begin{lemma}\label{lem:polykernel}
For each integer $m\ge1$,
\[
 \Ker(\D^m)=\Span_\C\{1,\omega,\ldots,\omega^{m-1}\}.
\]
\end{lemma}
\begin{proof}
The inclusion from right to left is the ordinary power rule. For the
reverse inclusion, induct on $m$. The case $m=1$ is the constant-field
property. If $\D^m y=0$, then $\D^{m-1}y=c\in\C$. Subtract
$c\omega^{m-1}/(m-1)!$ to reduce to the induction hypothesis.
\end{proof}

\begin{theorem}[Constant-coefficient solution-space classification]
\label{thm:constantpolynomial}
Let $0\ne P(T)\in\C[T]$, and let $m_\lambda$ be the multiplicity of
$\lambda$ as a root of $P$. The complete solution space in $\SC$ of
$P(\D)y=0$ is
\begin{equation}\label{eq:constantbasis}
 \bigoplus_{\substack{\lambda\in\R\\P(\lambda)=0}}
 e^{\lambda\omega}\,
 \Span_\C\{1,\omega,\ldots,\omega^{m_\lambda-1}\}.
\end{equation}
Its complex dimension is
$\sum_{\lambda\in\R,\,P(\lambda)=0}m_\lambda$.
For a nonzero constant polynomial this is the zero space.
\end{theorem}
\begin{proof}
For real $\lambda$ and arbitrary $v$,
\[
 (\D-\lambda)(e^{\lambda\omega}v)=e^{\lambda\omega}\D v.
\]
Iteration and \cref{lem:polykernel} identify the kernel of
$(\D-\lambda)^{m_\lambda}$ with the corresponding summand in
\eqref{eq:constantbasis}. For nonreal $\lambda$,
\cref{cor:constantlambda} says that $\D-\lambda$ is injective, so all
its positive powers are injective.

It remains to separate distinct roots. Factor $P$ into the pairwise
coprime powers $P_j(T)=(T-\lambda_j)^{m_j}$. The Chinese remainder
identities give polynomials $e_j(T)$ with
$\sum_j e_j\equiv1\pmod P$, such that $e_j$ is $1$ modulo $P_j$ and
$0$ modulo the other factors. For $y\in\Ker P(\D)$ put
$y_j=e_j(\D)y$. Then $\sum_jy_j=y$ and $P_j(\D)y_j=0$.
The same polynomial identities make the decomposition direct. The
nonreal-root summands vanish, leaving exactly \eqref{eq:constantbasis}.
Within each real-root summand the powers of $\omega$ are linearly
independent over $\C$, since $\omega$ is transcendental over the
ordinary constants. This gives the dimension.
\end{proof}

\begin{example}
For
\[
 P(T)=(T-2)^2(T+1)(T^2+1),
\]
the order is $5$, but the solution space inside $\SC$ has dimension $3$:
\[
 y=(c_0+c_1\omega)e^{2\omega}+c_2e^{-\omega},
 \qquad c_0,c_1,c_2\in\C.
\]
There are no additional sine and cosine terms associated with $T^2+1$.
In particular, $\D^2y+y=0$ has only the zero solution, even for complex
$y$.
\end{example}

\subsection{Constant matrices}
\begin{theorem}\label{thm:matrix}
Let $A\in M_n(\C)$. The solution space of $\D Y=AY$ in $(\SC)^n$ has
complex dimension equal to the sum of the algebraic multiplicities of the
\emph{real} eigenvalues of $A$. An invertible fundamental matrix in
$\GL_n(\SC)$ exists if and only if all eigenvalues of $A$ are real.
\end{theorem}
\begin{proof}
A similarity matrix over $\C$ is fixed by $\D$, so reduce to Jordan
form. On a Jordan block $\lambda I+N$ with real $\lambda$, a fundamental
matrix is
\[
 e^{\lambda\omega}\exp(\omega N)
 =e^{\lambda\omega}\sum_{k=0}^{m-1}\frac{\omega^kN^k}{k!}.
\]
This is a finite polynomial in the nilpotent matrix $N$, not an infinite
matrix exponential requiring convergence. Its inverse uses $-\omega N$.
On a nonreal-eigenvalue block, the last coordinate satisfies
$\D y_m=\lambda y_m$, so it is zero. Induction upward through the block
makes every coordinate zero. The dimension and invertibility assertions
follow block by block.
\end{proof}

\subsection{Global exponential and phase conventions cannot repair this}
\begin{corollary}[No globally compatible complex exponential]
\label{cor:noglobalexp}
There is no map $E:\SC\to\SC^\times$ satisfying
\begin{equation}\label{eq:forbiddenchain}
 \D(E(z))=E(z)\D z\qquad\text{for every }z\in\SC.
\end{equation}
In particular, a global complex exponential with this scalar-derivation
compatibility does not exist.
\end{corollary}
\begin{proof}
At $z=i\omega$, \eqref{eq:forbiddenchain} would produce the nonzero
solution $E(i\omega)$ of $\D y=iy$.
\end{proof}
Even the additive-to-multiplicative group law is unnecessary for this
contradiction. Similarly, no extension of finite $\cis$ to all real
surreals can satisfy $\D\cis(x)=i\D x\cis(x)$ everywhere.

This does not contradict algebraically defined global phase conventions
in the trigonometry report. It shows precisely which additional
compatibility they cannot possess. Nor does it deny that the ordinary
function $z\mapsto e^{iz}$ has a holomorphic germ or a finite-halo lift.
That germ is differentiated with respect to its argument. It is a
different assertion to evaluate a global exponential at $i\omega$ and
require the resulting \emph{number} to obey the selected field derivation.

\section{Where an oscillatory solution can live}\label{sec:PV}

The obstruction is internal to $\SC$ with its specified derivation. It
does not prevent constructing an abstract differential-field extension.
A completely explicit set-sized example makes the distinction precise.

\subsection{An algebraically closed base with no oscillatory element}
Let
\[
 K_0=\C((t^\Q)),\qquad
 \D\!\left(\sum_qc_qt^q\right)=\sum_q(-q)c_qt^{q+1}.
\]
This is the set-sized differential Hahn field of
\cref{ex:logintegral}, now complexified. Its underlying field is
algebraically closed and its differential constants are $\C$. It contains
$\omega=t^{-1}$, with $\D\omega=1$, and contains no nonzero solution of
$\D y=niy$ for $0\ne n\in\Z$, since it is a differential subfield of
$\SC$.

Adjoin a transcendental indeterminate $U$ and extend the derivation to
$K_0(U)$ by
\begin{equation}\label{eq:PVU}
 \D U=iU.
\end{equation}
The ordinary polynomial product rule and the quotient rule define the
extension uniquely. The formal generator $U$ is not a surreal number
endowed with the original $\D$.

\begin{theorem}[Explicit oscillatory Picard--Vessiot field]
\label{thm:PV}
The differential constants of $K_0(U)$ are exactly $\C$. The extension is
a Picard--Vessiot field for $\D y=iy$, and its differential automorphisms
over $K_0$ are precisely
\[
 U\longmapsto cU,\qquad c\in\C^\times.
\]
Thus its differential Galois group is the multiplicative group
$\mathbb G_m$ over $\C$.
\end{theorem}
\begin{proof}
Let $f\in K_0(U)^\times$ have $\D f=0$. Write
$f=cP/Q$ with $c\in K_0^\times$ and monic coprime polynomials
$P,Q\in K_0[U]$, of degrees $m,n$. After clearing denominators,
\[
 \frac{\D c}{c}PQ+(\D P)Q-P(\D Q)=0.
\]
Reduction modulo $P$ and modulo $Q$ gives $P\mid\D P$ and
$Q\mid\D Q$. The derivation does not increase degree in $U$. Its leading
coefficient on a monic degree-$m$ polynomial is $im$. Therefore
\[
 \D P=imP,\qquad\D Q=inQ,\qquad
 \frac{\D c}{c}=i(n-m).
\]
If $m\ne n$, the last equation is forbidden in $K_0$. Hence $m=n$ and
$c\in\C^\times$.

For a coefficient $a_j$ of $U^j$ in $P$, with $j<m$, coefficient
comparison gives
$\D a_j=i(m-j)a_j$. Again there is no nonzero such coefficient in
$K_0$. Thus $P=U^m$ and similarly $Q=U^n$. Coprimality and $m=n$ force
$m=n=0$. Hence $f\in\C^\times$, proving that no new constants arise.

The field is generated by the nonzero fundamental solution $U$ and has
no new constants, which is the Picard--Vessiot field condition for this
scalar equation. If $\sigma$ is a differential automorphism over $K_0$,
then $\D(\sigma(U)/U)=0$, so $\sigma(U)=cU$ for some
$c\in\C^\times$. Conversely every such scaling is an automorphism.
The scaling action realizes the stated algebraic group.
\end{proof}

\subsection{Why a formal exponential is not already in the base}
One might try to identify $U$ with
\[
 \sum_{n\ge0}\frac{(i\omega)^n}{n!}
 =\sum_{n\ge0}\frac{i^n}{n!}t^{-n}.
\]
Its exponent set $\{0,-1,-2,\ldots\}$ is not well ordered, so it is not a
Hahn series in $K_0$ or a summable surreal family. Choosing a larger
exponent group does not repair that support defect. The differential
extension is therefore not obtained simply by appending a missing value
to the list of permitted exponents.

Since $K_0$ is already algebraically closed, $U$ is necessarily
transcendental as a field element. It is nevertheless
\emph{differentially algebraic}, since it satisfies the nonzero
differential polynomial $\D U-iU=0$. Field-transcendental must not be
confused with differentially transcendental.

The preceding explicit extension proves, without invoking a general
existence theorem, that algebraic closedness and a rich scalar calculus do
not make the original field differentially closed. For the full $\SC$,
one can similarly formulate an extension symbolically, but the set-sized
construction above avoids any need to form a field of class-sized
rational expressions.

\section{The analytic coefficient interface}\label{sec:analytic}

This section connects the new scalar structure to the common-domain
Hahn calculus in the repository. We do not transport statements between
arbitrary formal germs and fixed-domain holomorphic families.

\subsection{Differentiating an analytic Hahn family}
Let $V\subseteq\C$ be an ordinary nonempty open set. Define
\[
 \HH_\Gamma(V)=\OO(V)((t^\Gamma))
\]
to consist of formal families
\begin{equation}\label{eq:analyticfamily}
 F(z)=\sum_{\gamma\in S}f_\gamma(z)t^\gamma,
 \qquad f_\gamma\in\OO(V),
\end{equation}
with a common well-ordered support $S\subseteq\Gamma$.
The symbol $\OO(V)$ here denotes ordinary holomorphic functions; it is
not the finite-element valuation ring $\OO$ without an argument.
The analytic derivative is
$\partial_zF=\sum_\gamma f'_\gamma(z)t^\gamma$.

For each $\gamma$ write
\begin{equation}\label{eq:monomialmatrix}
 \D t^\gamma=\sum_\delta d_{\delta\gamma}t^\delta,
 \qquad d_{\delta\gamma}\in\R.
\end{equation}
These scalars form a support-controlled matrix, not a matrix indexed by
all surreals as a set. Only set-sized columns are used at a time.

\begin{lemma}[Finite incoming contributions]\label{lem:finiteincoming}
For every well-ordered set $S$ of exponents, the union
\[
 T_S=\bigcup_{\gamma\in S}\supp\D t^\gamma
\]
is well ordered, and for each $\delta$ only finitely many
$\gamma\in S$ have $d_{\delta\gamma}\ne0$.
\end{lemma}
\begin{proof}
The family $(t^\gamma)_{\gamma\in S}$ is summable. Apply strong
additivity, including its assertion that the derivative family is
summable. The two conclusions are exactly the two clauses of
\cref{def:summable}.
\end{proof}

\begin{theorem}[Support-controlled common-domain interface]
\label{thm:analyticinterface}
Let $\Gamma$ be a differentially stable set-sized exponent group. The
rule
\begin{equation}\label{eq:coefficientderivation}
 \D_HF
 =\sum_{\delta\in T_S}
       \left(\sum_{\gamma\in S}d_{\delta\gamma}f_\gamma\right)t^\delta
\end{equation}
defines a $\C$-linear derivation
$\D_H:\HH_\Gamma(V)\to\HH_\Gamma(V)$.
Every inner sum is finite. In particular, differentiating the scalar
coefficients does not shrink or destroy the common ordinary domain $V$.
Moreover,
\begin{equation}\label{eq:commuteanalytic}
 \D_H\partial_zF=\partial_z\D_HF.
\end{equation}
For an arbitrary initial $\Gamma$, the same assertions hold after
passing to its differential Hahn hull.
\end{theorem}
\begin{proof}
\Cref{lem:finiteincoming} proves the well-ordered support and finiteness
of each inner sum. Each new coefficient is a finite real-linear
combination of holomorphic functions on the \emph{same} $V$, and so is
holomorphic there. The supports lie in $\Gamma$ by stability.

Linearity is immediate. To verify the product rule, expand the Hahn
product of two families and apply
$\D(t^{\gamma+\eta})=t^\gamma\D t^\eta+t^\eta\D t^\gamma$.
The standard product summability and strong additivity make the expanded
families summable; coefficient comparison then uses only finite sums.
Finally ordinary differentiation is complex-linear and commutes with each
finite inner sum in \eqref{eq:coefficientderivation}. This gives
\eqref{eq:commuteanalytic}. The hull theorem supplies the last assertion.
\end{proof}

\subsection{Linear analytic operations commute with the scalar derivation}
The proof uses only the linearity of $f\mapsto f'$. It therefore gives a
more general useful statement.

\begin{corollary}[Coefficientwise linear naturality]\label{cor:naturality}
Any fixed $\C$-linear operation on ordinary coefficient functions commutes
with $\D_H$, wherever that operation is defined coefficientwise. In
particular, for a fixed ordinary piecewise smooth contour $C\subseteq V$,
\begin{equation}\label{eq:contourcommute}
 \D\left(\int_C F(z)\dd z\right)
 =\int_C\D_HF(z)\dd z,
 \qquad
 \int_C F\dd z:=\sum_\gamma
      \left(\int_C f_\gamma\dd z\right)t^\gamma.
\end{equation}
The analogous identity holds for Laurent-coefficient residues at a fixed
ordinary center whenever all coefficient functions belong to the stated
punctured-domain class.
\end{corollary}
\begin{proof}
Apply the fixed linear operation to the finite inner sums in
\eqref{eq:coefficientderivation}. Both orders of application have
coefficient $\sum_\gamma d_{\delta\gamma}L(f_\gamma)$ at $\delta$.
\end{proof}
The contour is fixed and ordinary. A moving contour, moving endpoints,
or a residue with a different center convention requires additional
terms or a different statement; none is hidden in this corollary.

For example, on a punctured ordinary disk about $a\in\C$,
$F(z)=t/(z-a)$ has contour integral $2\pi it$ around a positive simple
loop. Both sides of \eqref{eq:contourcommute} are $-2\pi it^2$. This is
an exact Hahn identity.

\subsection{Moving evaluation and the correct total derivative}
For $c\in V$ and $h\in\mm_\C$, define the finite-halo evaluation of
\eqref{eq:analyticfamily} by
\begin{equation}\label{eq:haloevaluation}
 F^\#(c+h)=
 \sum_{\gamma\in S}\sum_{n\ge0}
     \frac{f_\gamma^{(n)}(c)}{n!}\,t^\gamma h^n.
\end{equation}
The family is summable: its support is controlled by the sum of the
well-ordered $S$ and the positive-support monoid generated by $h$.
Finiteness of representations in the Hahn product and the
Hahn--Neumann lemma give coefficientwise finiteness. No uniform bound on
the ordinary Taylor coefficients is required for this strong-summation
statement.

\begin{theorem}[Total-derivative chain rule]\label{thm:totalchain}
With the common-domain hypotheses above,
\begin{equation}\label{eq:totalchain}
 \D\bigl(F^\#(c+h)\bigr)
 = (\D_HF)^\#(c+h)
   +\D h\,(\partial_zF)^\#(c+h).
\end{equation}
All terms may be placed in one set-sized differentially stable Hahn
workspace after adjoining the support of $h$ and taking the hull.
\end{theorem}
\begin{proof}
Differentiate the summable family in \eqref{eq:haloevaluation}. The
ordinary complex Taylor coefficients have derivative zero, so the
product rule splits each term into
\[
 \frac{f_\gamma^{(n)}(c)}{n!}
 \bigl((\D t^\gamma)h^n
      +t^\gamma n h^{n-1}\D h\bigr),
\]
with the second term omitted for $n=0$. The first family is summable by
\cref{lem:finiteincoming} followed by the same positive-support Taylor
argument. The second is the corresponding Taylor family for
$\partial_zF$, multiplied by the fixed scalar $\D h$. Thus the split
and rearrangement are legitimate, and they give the two terms of
\eqref{eq:totalchain}. The workspace assertion follows from
\cref{thm:hull}.
\end{proof}

\begin{example}\label{ex:totalchain}
Take $F(z)=t^2e^z$ and evaluate at $z=t$. The three relevant expressions
are
\[
 \D_HF=-2t^3e^z,\qquad
 \partial_zF=t^2e^z,\qquad
 \D(t^2\Expzero(t))=(-2t^3-t^4)\Expzero(t).
\]
The extra $-t^4\Expzero(t)$ is the moving-argument term. Differentiating
only the scalar prefactor or only the analytic argument would miss part
of the answer.
\end{example}

\subsection{Formal power series need their own admissibility condition}
For a general $F(Z)=\sum_{n\ge0}c_nZ^n\in K[[Z]]$, formal coefficient
and variable derivatives commute:
\[
 \D_cF=\sum_n(\D c_n)Z^n,\qquad
 \partial_ZF=\sum_{n\ge1}n c_nZ^{n-1},\qquad
 \D_c\partial_ZF=\partial_Z\D_cF.
\]
Evaluation at $a$ requires summability. If the families defining
$F(a)$, $(\D_cF)(a)$, and $(\partial_ZF)(a)$ are summable, then the same
calculation proves
\[
 \D(F(a))=(\D_cF)(a)+\D a\,(\partial_ZF)(a).
\]
One must not apply this formula to an arbitrary unevaluable formal germ.
The common-domain finite-halo theorem supplies its summability proof;
an unrestricted formal coefficient sequence does not. This is exactly
the kind of interface boundary emphasized in the repository's analysis
and foundations reports \cite{RepoAnalysis,RepoFound}.

\section{Transseries, composition, and the limits of transfer}
\label{sec:transfer}

Aschenbrenner, van den Dries, and van der Hoeven prove that the natural
embedding of the ordered differential field of logarithmic-exponential
transseries into $(\No,\D)$ is elementary, and that Hardy fields embed
in the appropriate ordered differential sense \cite{ADH}. This places the
selected derivation in a powerful nonoscillatory asymptotic setting. It is
not a theorem that every algebraic differential equation has a solution.
The explicit obstruction $\D y=iy$, $y\ne0$, already shows why
complexification must not be called a differentially closed field.

Berarducci and Mantova construct composition and a compatible Taylor
calculus on the omega-series field. Their global incompatibility theorem
says that the particular simplest derivation used here cannot be
compatible with a composition on all of $\No$ under their stated
compatibility axioms \cite{BMComposition}. Thus one must distinguish
composition on its proved domain from a hypothetical functional
interpretation of every surreal number. No assertion here settles the
broader problem of alternative global composition structures.

Several tempting inferences can now be rejected by concrete examples.
Algebraic closedness does not imply logarithmic-derivative surjectivity:
$i$ is missing from the image. Surjectivity of $\D$ does not imply
closure under all exponential integrals: $i\omega$ is a primitive of $i$,
but there is no compatible exponential of it in $\SC$. Differential
stability of a Hahn workspace does not imply primitive closure:
$\log\omega$ is missing from $H_\Q$. Infinitesimal imaginary part is
necessary but not sufficient for a nonzero rank-one solution:
$i/\omega$ is excluded whereas $i/\omega^2$ is included.

There is also no contradiction between the missing harmonic oscillator
and the existence of finite-angle sine and cosine. If
$s=\sin(\varepsilon)$ and $c=\cos(\varepsilon)$ for an infinitesimal
$\varepsilon$, then the proved identities are
\[
 \D s=c\D\varepsilon,\qquad
 \D c=-s\D\varepsilon.
\]
The factor $\D\varepsilon$ is indispensable. Removing it would replace
an infinitesimal phase speed by the forbidden constant speed $1$.

\section{Computation and formalization contracts}\label{sec:computation}

\subsection{What the accompanying program actually implements}
The package includes an exact Python/SymPy verifier and a small decision
routine for the rational criterion of \cref{thm:rational}. The routine
accepts $b(X)\in\Q(X)$, cancels common factors, and reports whether the
degree gap is at least two. Zero is handled separately. Inputs involving
other functions or nonrational coefficients are rejected as outside the
implemented domain, not classified as mathematical nonexistence.

The verification program checks finite rational identities, truncated
small-exponential equations and their exact residuals, the algebraic
solution of \cref{ex:algebraicphase}, chain-rule examples, a finite Jordan
fundamental matrix, and the displayed constant-coefficient basis.
The symbol $X$ in these checks is an ordinary formal real variable. A
successful symbolic derivative check validates the displayed finite
identity; it does not implement the Berarducci--Mantova derivation on an
arbitrary surreal normal form.

For the truncated exponential
\[
 E_N(t)=\sum_{n=0}^N\frac{(-i)^n}{n!}t^n,
 \qquad\D=-t^2\frac{\dd}{\dd t},
\]
the exact residual is
\begin{equation}\label{eq:truncresidual}
 \D E_N-it^2E_N
 =-i\frac{(-i)^N}{N!}t^{N+2}.
\end{equation}
This identity, not numerical convergence, is the finite verification target.
The infinite identity follows separately from strong summability.

\subsection{A scalar derivative needs a monomial support contract}
A useful extension to a symbolic Hahn representation should supply both
$\D(t^\gamma)$ and its support certificate. A function returning a formal
expression without an admissible support does not yet implement a strong
Hahn derivation. The hull theorem supplies a semantic closure guarantee,
but its countably many stages do not constitute a terminating algorithm
for every input representation.

For groups contained in $\R$, \cref{prop:realgroup} gives an unusually
simple effective closure: adjoin exponent $1$ and its rational multiples.
For genuinely surreal exponent groups, a derivative may introduce new
logarithmic scales. The universal power shortcut warned against after
\cref{prop:calculus} must not be used to hide that enlargement.

The existing computer-algebra report separates exact denotation,
coefficient access, and decision procedures \cite{RepoCAS}. The same
separation applies here. Exact notation $\Ip(b)$ is not an algorithm for
its leading purely infinite term. An obstruction certificate should be a
proved nonzero value of $\pinf\Ip(b)$, or a theorem-specific witness such
as the degree test, rather than a timeout in an integration routine.

\subsection{A proof-assistant interface, without an implementation claim}
A future formalization should distinguish a scalar derivation on a
set-sized coefficient field, a strongly additive Hahn lift of it, formal
variable differentiation, and an analytic evaluation map. The most useful
new obligation is the implication
\[
 \text{monomial support closure}
 \quad\Longrightarrow\quad
 \text{closure of the full small-support Hahn type under }\D.
\]
This is the proof content of \cref{lem:monomialcriterion,thm:hull}.
It complements, rather than duplicates, generic Hahn convolution.

A second obligation is \cref{lem:finiteincoming}: for each admissible
input support, each output coefficient receives only finitely many
monomial-derivative contributions. Once available, the common-domain
analytic coefficient interface is finite linear algebra at every output
exponent. A third is the represented quotient $\JJ$ in
\cref{rem:quotient}, so that rank-one classification does not accidentally
require a type whose elements are proper classes.

No Lean files are delivered or asserted to compile. In particular, this
article does not claim that the concrete surreal carrier has already been
connected to every required Hahn, real-closed-field, exponential, or
derivation interface. The elementary classification can first be
formalized over a differential ordered field satisfying the explicitly
listed axioms, with the concrete surreal instantiation left as a separate
proved dependency.

\section{Conclusions and precise remaining tasks}\label{sec:conclusion}

The central missing connection is now explicit. A scalar derivation can
be made internal to a set-sized Hahn workspace, and it can act on a
common-domain analytic Hahn family without destroying the common domain.
Finite phase is not merely a geometric convenience: it exactly controls
which imaginary logarithmic derivatives can occur. The purely infinite
part of an imaginary primitive is a complete rank-one obstruction.

The higher-rank situation remains a separate problem. Constant matrices
are fully handled by \cref{thm:matrix}; arbitrary matrices with surreal
entries are not. Rank-one gauge classification suggests studying how
infinite phase obstructions interact with invariant subspaces and changes
of basis, but no general matrix normal form follows from the scalar
result. Likewise, an effective description of differential hulls for
large symbolic exponent languages needs a separately specified
representation and monomial-derivative procedure.

For the repository, a concrete next mathematical task is to combine the
coefficient derivation with its existing finite-deformation and residue
interfaces, always fixing whether the centers and contours move. The
fixed-domain, fixed-contour compatibility proved here is the appropriate
starting theorem. It is not permission to differentiate every residue
convention as though the center were constant.

The resulting picture is coherent rather than paradoxical: $\SC$ is
algebraically closed; $\D$ is surjective; real exponential integrals
exist; finite angular motion is exact; and sustained nonzero constant
angular velocity is nevertheless absent. These statements concern
different closure properties, and the obstruction sequence
\eqref{eq:exactsequence} records their precise boundary.

\clearpage
\appendix
\section{Dependency and verification ledger}\label{app:ledger}

\begin{longtable}{L{0.30\textwidth}L{0.33\textwidth}L{0.28\textwidth}}
\toprule
Claim & Mathematical input & Verification status\\
\midrule
\endhead
Existence and surjectivity of $\D$ & Berarducci--Mantova theorem & Imported,
not reproved or machine-checked\\
Hahn workspace construction & Set-supported normal forms; strong additivity &
Full proof in \cref{sec:hull}\\
Finite phase & Positive-support formal exponential and logarithm & Full
proof in \cref{sec:phase}; finite series checks supplement it\\
Finite-primitive subgroup & $H$-field sign, normalization, primitives & Full
proof in \cref{sec:bounded}\\
Logarithmic-derivative image and gauge classes & Polar decomposition and
real exponential integrals & Full proof in \cref{sec:logder}\\
Rational decision criterion & Leading Laurent exponent; convexity of $\BB$ &
Full proof; exact restricted implementation\\
Constant-coefficient equations & Constant field, no oscillation, polynomial
Bezout identities & Full proof; displayed examples symbolically checked\\
Oscillatory extension & Rational-function divisibility and no constant
imaginary logarithmic derivative & Full no-new-constants proof\\
Analytic interface & Finite incoming derivative contributions; common
ordinary domain & Full proof; worked examples checked\\
Proof-assistant instantiation & Concrete surreal/Hahn bridge and chosen
scalar derivation & Not implemented in this package\\
\bottomrule
\end{longtable}

The generated verification record is \path{data/verification.json}.
It contains the actual executed check count, software versions, and the
statement of what was tested. \path{data/build_report.json} records the
PDF build and layout checks. Neither file certifies a mathematical theorem
beyond the finite identities explicitly executed. The paper's general
proofs remain human-readable mathematical arguments, not kernel-checked
proof objects.

\section{Build and package contents}\label{app:build}
The article is standalone: the bibliography is embedded and no external
figures or bibliography processor are required. From the package directory,
run
\begin{lstlisting}
latexmk -pdf -interaction=nonstopmode -halt-on-error article.tex
python code/verify_examples.py
\end{lstlisting}
Alternatively, run \path{build.sh}; it uses three pdfLaTeX passes if
\texttt{latexmk} is unavailable. The verification script requires Python
and SymPy; it uses exact arithmetic rather than floating-point simulation.
The script can also be imported to use the restricted rational-phase
classification routine.

The repository audit is recorded separately in \path{REPOSITORY_AUDIT.md}
with the inspected commit and document paths. It documents a targeted gap
review and does not claim to certify the existing reports. Only the new
article and its supporting artifacts are packaged; the repository was not
modified.

\begin{thebibliography}{99}

\bibitem{BM}
A. Berarducci and V. Mantova,
\emph{Surreal numbers, derivations and transseries},
Journal of the European Mathematical Society \textbf{20} (2018), 339--390.
DOI: \href{https://doi.org/10.4171/JEMS/769}{10.4171/JEMS/769}.
Public preprint: \url{https://arxiv.org/abs/1503.00315}; the theorem
numbering A and B used here is visible in version 3.

\bibitem{ADH}
M. Aschenbrenner, L. van den Dries, and J. van der Hoeven,
\emph{The surreal numbers as a universal $H$-field},
arXiv:1512.02267, version 3 (2016).
\url{https://arxiv.org/abs/1512.02267}.
Theorems 1--3 and the Hahn-field preliminaries are the relevant passages.

\bibitem{BMComposition}
A. Berarducci and V. Mantova,
\emph{Transseries as germs of surreal functions},
Transactions of the American Mathematical Society \textbf{371} (2019),
3549--3592.
\url{https://arxiv.org/abs/1703.01995}.
See the composition, Taylor, and incompatibility theorems (Theorems 6.3,
7.14, and 8.4 in the consulted preprint).

\bibitem{Neumann}
B. H. Neumann,
\emph{On ordered division rings},
Transactions of the American Mathematical Society \textbf{66} (1949),
202--252.
DOI: \href{https://doi.org/10.1090/S0002-9947-1949-0032593-5}
{10.1090/S0002-9947-1949-0032593-5}.
\url{https://www.jstor.org/stable/1990552}.

\bibitem{RepoMap}
V. Reshetnikov, \texttt{Surreal} repository,
\emph{The surreal and surcomplex reports},
\path{docs/README.md} and \path{docs/manifest.tex},
commit \texttt{e260237db9b71da8b74a0c13c8e6355119091100},
inspected September 21, 2026.
\url{https://github.com/VladimirReshetnikov/Surreal/tree/e260237db9b71da8b74a0c13c8e6355119091100/docs}.

\bibitem{RepoTrig}
\texttt{Surreal} repository,
\emph{Trigonometry on the Surcomplex Plane}, specialist README,
\path{docs/surcomplex/trigonometry/README.md}, same inspected commit.
\url{https://github.com/VladimirReshetnikov/Surreal/blob/e260237db9b71da8b74a0c13c8e6355119091100/docs/surcomplex/trigonometry/README.md}.

\bibitem{RepoAnalysis}
\texttt{Surreal} repository,
\emph{Surcomplex analysis: holomorphic functions on $K=\No[i]$},
specialist README,
\path{docs/surcomplex/analysis/README.md}, same inspected commit.
\url{https://github.com/VladimirReshetnikov/Surreal/blob/e260237db9b71da8b74a0c13c8e6355119091100/docs/surcomplex/analysis/README.md}.

\bibitem{RepoFound}
\texttt{Surreal} repository,
\emph{Foundations for surreal and surcomplex mathematics},
specialist README,
\path{docs/foundations-and-computation/foundations/README.md},
same inspected commit.
\url{https://github.com/VladimirReshetnikov/Surreal/blob/e260237db9b71da8b74a0c13c8e6355119091100/docs/foundations-and-computation/foundations/README.md}.

\bibitem{RepoCAS}
\texttt{Surreal} repository,
\emph{Surreal and Surcomplex Numbers in Symbolic Computer Algebra},
specialist README,
\path{docs/foundations-and-computation/computer-algebra/README.md},
same inspected commit.
\url{https://github.com/VladimirReshetnikov/Surreal/blob/e260237db9b71da8b74a0c13c8e6355119091100/docs/foundations-and-computation/computer-algebra/README.md}.

\end{thebibliography}
\end{document}
